\newglossaryentry{mpc-term}{
    name={Многосторонние конфиденциальные вычисления},
    description={совместное вычисление заданной функции, при котором сведения, доступные нарушителю, ограничены его входами, разрешёнными выходами и публичными параметрами в принятой модели безопасности}
}

\newglossaryentry{secret-sharing}{
    name={Разделение секрета},
    description={представление значения набором долей, позволяющее восстановить значение при объединении долей уполномоченного множества участников}
}

\newglossaryentry{replicated}{
    name={Реплицированное разделение},
    description={способ размещения долей секрета, при котором одну компоненту хранят несколько сторон. В рассматриваемой схеме каждый из трёх серверов хранит две из трёх аддитивных компонент}
}

\newglossaryentry{semi-honest}{
    name={Пассивный (получестный) нарушитель},
    description={участник, который следует протоколу и использует доступные ему доли, случайные значения и сообщения для получения сведений о чужих входах}
}

\newglossaryentry{static}{
    name={Статический нарушитель},
    description={нарушитель, множество скомпрометированных сторон которого фиксируется до начала исполнения протокола}
}

\newglossaryentry{threshold}{
    name={Порог компрометации},
    description={максимальное число совместно наблюдаемых нарушителем серверов, при котором сохраняется рассматриваемая гарантия. В принятой модели пассивного нарушителя порог равен одному серверу из трёх}
}

\newglossaryentry{decision-tree}{
    name={Решающее дерево},
    description={модель классификации, в которой внутренние узлы задают условия выбора ветви, а листья содержат метки классов}
}

\newglossaryentry{private-inference}{
    name={Конфиденциальный вывод},
    description={вычисление прогноза по заданной модели с защитой закрытых входов и промежуточных значений и с выдачей результата назначенному получателю}
}

\newglossaryentry{fixed-point}{
    name={Кодирование с фиксированной точкой},
    description={представление числа целым кодом, полученным умножением на общий масштаб и округлением. Для признаков и порогов в работе применяется масштаб $2^f$}
}

\newglossaryentry{reference}{
    name={Целочисленный эталон},
    description={открытый обход дерева по закодированным признакам и порогам с тем же правилом сравнения, которое используется при конфиденциальном выводе}
}

\newacronym{mkv}{МКВ}{многосторонние конфиденциальные вычисления}
\newacronym{aby}{ABY3}{трёхсторонний протокол с арифметическим и двоичным представлениями и схемами Яо (Arithmetic, Binary, Yao)}
\newacronym{tfe}{TF Encrypted}{платформа конфиденциальных вычислений на основе TensorFlow}

\newglossaryentry{sym-tree}{
    type=\acronymtype, name={$T$}, sort={tree},
    description={решающее дерево}
}

\newglossaryentry{sym-input}{
    type=\acronymtype, name={$x$}, sort={input},
    description={вектор признаков объекта}
}

\newglossaryentry{sym-depth}{
    type=\acronymtype, name={$d$}, sort={depth},
    description={глубина дерева в рёбрах}
}

\newglossaryentry{sym-nodes}{
    type=\acronymtype, name={$m$}, sort={nodes},
    description={число внутренних узлов дерева}
}

\newglossaryentry{sym-batch}{
    type=\acronymtype, name={$B$}, sort={batch},
    description={число объектов в пакете}
}

\newglossaryentry{sym-precision}{
    type=\acronymtype, name={$f$}, sort={precision},
    description={число дробных разрядов}
}

\newglossaryentry{sym-impl}{
    type=\acronymtype, name={$e_{\mathrm{impl}}$}, sort={impl},
    description={доля расхождений между конфиденциальным выводом и целочисленным эталоном}
}

\newglossaryentry{sym-quant}{
    type=\acronymtype, name={$e_{\mathrm{quant}}$}, sort={quant},
    description={доля расхождений между целочисленным эталоном и открытым выводом по исходному дереву}
}
